\section{Domain of validity: exact ingredients and approximations}
\label{sec:validity}

Before exploiting the results of the previous sections, it is useful to state precisely what they
contain and which approximations, with respect to a calculation in full QCD, have been made.

\paragraph{What is treated exactly.} Within the perturbative order and the high-energy approximation
specified below, the following ingredients are exact.
\begin{itemize}
\item \emph{Kinematics of the projectile.} The dependence on the quark mass $\mf$, on the virtuality
$Q^2$, on the longitudinal fractions $\zeta$, $\eta$, $\xi$ and on the three transverse momenta is
complete. No soft-gluon, collinear, small-mass, heavy-quark or back-to-back approximation has been made.
In particular the results hold for $Q^2\to0$ and for arbitrary ratios $\mf^2/Q^2$.
\item \emph{The partonic amplitude.} All the tree-level graphs for $\gamma_T^*\to\qqg$ are included: both
emitters, the regular and the instantaneous terms, and all helicity and polarisation states. The
two-component expressions that we use have been checked against the amplitude built with
four-component Dirac spinors (appendix~\ref{app:checks}).
\item \emph{Multiple scattering.} The interaction with the target is included to all orders in the
strong field $gA^-\sim1$ through the Wilson lines; the result therefore resums all the ``higher-twist''
corrections in powers of $Q_s^2/k_\perp^2$, where $Q_s$ is the saturation momentum of the target.
\item \emph{Colour.} The cross-sections \eqref{eq:xs-qbar}, \eqref{eq:xs-q} and \eqref{eq:xs-int} with
the colour structure \eqref{eq:S6}--\eqref{eq:W} are valid for any number of colours.
\end{itemize}

\paragraph{Approximations.} Our results are \emph{not} full QCD results. They rely on the following
approximations.
\begin{enumerate}
\item \emph{Perturbative order.} The cross-section is of order $\aem\as$, i.e.\ it is the leading-order
(tree-level) result for the three-parton final state. Seen as a contribution to the production of a
heavy-quark pair, it provides only the real part of the NLO corrections; the virtual corrections,
which involve the one-loop $\qqbar$ component of the photon wave function with massive
quarks~\cite{Beuf:2021qqa,Beuf:2021srj,Beuf:2022ndu}, are not included. The coupling is fixed.
\item \emph{Eikonal approximation.} The scattering off the target is treated at leading power in the
collision energy: the target is a shock wave of vanishing width, and each parton keeps its transverse
position, its longitudinal momentum and its helicity while crossing it. Corrections suppressed by
powers of the energy (finite width of the target, transverse and ``plus'' components of the target
field, quark exchanges with the target, spin-dependent
interactions)~\cite{Altinoluk:2014oxa,Kovchegov:2015pbl} are neglected. This requires the coherence
time of the three-parton fluctuation to be much larger than the longitudinal extent of the target, i.e.
\begin{equation}
x_{\rm eff}\equiv\frac{Q^2+M_{\qqg}^2}{W^2+Q^2}\ll1\,,\qquad
M_{\qqg}^2=\frac{\kk_1^2+\mf^2}{\zeta}+\frac{\kk_2^2+\mf^2}{\eta}+\frac{\kk_3^2}{\xi}
\quad(\kk_1+\kk_2+\kk_3=0)\,,
\end{equation}
with $W$ the photon--target centre-of-mass energy. For heavy quarks $M_{\qqg}^2\geq4\mf^2$, which
makes this condition more restrictive than for light quarks.
\item \emph{CGC description of the target.} The target is described by a classical colour field with a
statistical weight, and the observable is obtained by averaging products of Wilson lines. The
leading logarithms of the energy are resummed by the B--JIMWLK evolution of these averages. At the
accuracy of the present calculation the rapidity scale at which the Wilson lines are evaluated is not
fixed; it is fixed only together with the NLO corrections, whose soft-gluon part is discussed in
section~\ref{sec:soft}.
\item \emph{Partonic final state.} The cross-section is given for on-shell partons. For jets one must
apply a jet definition and for hadrons one must convolute with fragmentation functions. The partonic
cross-section is singular for $\xi\to0$ (soft gluon) and, for massless quarks, when the gluon is
collinear to its emitter; with massive quarks the collinear singularity is replaced by logarithms of
$k_\perp^2/\mf^2$ (section~\ref{sec:coll}).
\item \emph{Photon polarisation.} Only the transverse polarisations of the virtual photon are
considered. The longitudinal polarisation, and the interference between the two in the leptonic
cross-section, must be added to describe electroproduction at $Q^2\neq0$; in the photoproduction limit
the transverse photon is all one needs.
\end{enumerate}

\paragraph{Additional, optional approximations.} Some of the formulas involve further simplifications
of the colour correlators, which are not needed for the validity of the general result: the large-$\Nc$
limit in eqs.~\eqref{eq:W-largeN}--\eqref{eq:W-int}, which neglects corrections of relative order
$1/\Nc^2$; the Gaussian approximation \eqref{eq:Gauss} for the quadrupole; and, in
section~\ref{sec:mom}, the restriction to the coherent part of the cross-section for a target which is
homogeneous in the transverse plane.
